\section*{Chapitre \ref{ch:mx:build-an-execution-algorithm} --- Construction~: un algorithme d'exécution}

\begin{solution}{exo:mx:build-an-execution-algorithm:1}
5\,000 actions (une droite). La tolérance vaut $20/600\times10\,000=333$ actions, et l'algorithme croise pour ce qu'il a de retard au-delà de cette tolérance, en
lots entiers~: il croise dès qu'il est en retard d'au moins 433 actions, soit 500 actions, puisque ses exécutions arrivent par lots.
\end{solution}

\begin{solution}{exo:mx:build-an-execution-algorithm:2}
Les autres ont négocié $40\,000-8\,000=32\,000$ actions~; 30~\% du total autorise $0.3/0.7\times32\,000=13\,714$ actions pour lui, soit 5\,714 actions de plus
au regard du volume du marché jusqu'ici.
\end{solution}

\begin{solution}{exo:mx:build-an-execution-algorithm:3}
$1.25-0.70\times2-0.69=-0.84$ points de base (une perte~; la moyenne de la cellule vaut $-0.63$). Seuil de rentabilité à faible \omterm{def:mx:the-almgren-chriss-framework:urgency}{urgence}~: $\log_2(Q/4000)=1.25/0.70=1.78$, soit
$Q\approx13\,700$ actions.
\end{solution}

\begin{solution}{exo:mx:build-an-execution-algorithm:4}
Il prend la contrepartie des ordres de vente qui auraient sans cela frappé le reste de la file acheteuse et fait baisser le cours~: le cours reste plus élevé
qu'il ne l'aurait été sans l'ordre. Le prix milieu contrefactuel dans \texttt{firm.tca} mesure cet effet directement~; le rattrapage par croisement et l'avance des exécutions
s'y ajoutent.
\end{solution}

\begin{solution}{exo:mx:build-an-execution-algorithm:5}
Rattraper le retard reviendrait à croiser le spread pour une minute d'échéancier dès la réouverture du titre, alors que le carnet est mince et que les prix ne sont pas stabilisés~;
décaler l'échéancier conserve la participation prévue, au prix d'une fin plus tardive (ici une minute après l'échéance~; une échéance ferme imposerait au contraire de répartir les
actions manquantes sur le temps restant).
\end{solution}

\begin{solution}{exo:mx:build-an-execution-algorithm:6}
Il faut considérer les exécutions, et non l'ordre, comme la position~: réconcilier la drop copy après l'arrêt, signaler l'exécution tardive au trader et faire en sorte que le risque voie la position réelle~;
un coupe-circuit qui suppose que l'ordre s'est arrêté à l'instant de l'arrêt en sous-estime la taille.
\end{solution}

\begin{solution}{exo:mx:build-an-execution-algorithm:7}
Avant frais, l'économie vaut en moyenne $-0.47$ point de base~: les exécutions de l'algorithme sont moins bonnes de 0.47 point de base, mais il paie environ 0.32 de moins en frais
(il touche des rabais sur ses exécutions passives, $-0.02$ tick par action, là où le TWAP paie 0.30 pour prendre de la liquidité), si bien qu'au total il coûte seulement 0.15 de plus. Une comparaison avant
frais surestimerait la perte d'un facteur trois.
\end{solution}

\begin{solution}{exo:mx:build-an-execution-algorithm:8}
Les ordres qu'il a reçus ne sont pas les ordres du TWAP~: sans assignation aléatoire (ou une comparaison appariée sur les mêmes ordres), l'écart mêle
l'algorithme à la difficulté des ordres (chapitre~19)~; 200 ordres réels avec une dispersion appariée supérieure à 1.3 point de base donnent au mieux
un intervalle d'environ $\pm0.2$~; et une économie qui dépend de la taille ne dit rien des ordres d'un autre client. Il faut randomiser, ajuster selon la difficulté, présenter
l'intervalle et préciser les tailles auxquelles il s'applique.
\end{solution}

\begin{solution}{pb:mx:build-an-execution-algorithm:1}
\textbf{1.} Nouvel, en cours, suspendu, terminé, expiré, arrêté~; heure de début, arrêt ou pause et reprise, exécution, échéance, arrêt d'\omterm{def:mx:the-almgren-chriss-framework:urgency}{urgence}.
\textbf{2.} Heure, état, événement et motif de chaque transition et de chaque refus~: pour le desk, la conformité et le développeur.
\textbf{3.} Voir les définitions.
\textbf{4.} Chaque seconde~: les actions dues à présent par l'échéancier et celles dues dans 20 secondes, bornées par la bande~; un croisement pour ce qui est en retard au-delà de la tolérance,
balayé sur les plateformes~; un ordre passif au meilleur cours proche pour le reste de l'échéancier à venir.
\textbf{5.} Il laisse à l'ordre passif le temps de s'exécuter avant de croiser~; à l'échéance, le reste est croisé dans la limite, et ce qui n'est pas
exécuté expire.
\textbf{6.} Taille maximale d'un \omterm{def:mx:the-almgren-chriss-framework:parent}{ordre enfant}, un collar autour du médian, et jamais plus d'exécutions en cours ou réalisées que l'\omterm{def:mx:the-almgren-chriss-framework:parent}{ordre parent}.
\textbf{7.} Des contrôles raisonnablement conçus pour empêcher les ordres qui dépassent les seuils de crédit ou de capital, ainsi que les ordres erronés qui excèdent les paramètres de prix ou de taille
ou les ordres faisant double emploi.
\textbf{8.} Il annule ses ordres et arrête l'horloge, puis décale l'échéancier de la durée de la pause, afin d'éviter de se précipiter dans un marché en réouverture.
\textbf{9.} Il a cessé de croiser pendant 213 secondes et a continué d'acheter passivement à la limite~; 3.79 points de base contre 4.31 sans lui.
\textbf{10.} Un \omterm{def:mx:the-almgren-chriss-framework:parent}{ordre enfant} était en vol au moment où l'arrêt a été envoyé.
\textbf{11.} Quatre tailles, de 4\,000 à 32\,000 actions, sur dix minutes, deux niveaux d'\omterm{def:mx:the-almgren-chriss-framework:urgency}{urgence}, 25 graines~; le TWAP envoie un \omterm{def:mx:the-limit-order-book:orders}{ordre au marché} toutes les dix secondes pour ce que la droite de l'échéancier
fixe comme dû.
\textbf{12.} Pour que les deux algorithmes affrontent le même flux d'ordres et que l'écart mesure les algorithmes, et non les journées.
\textbf{13.} Le spread, l'impact (par rapport à une séance contrefactuelle), le timing, l'opportunité et les frais.
\textbf{14.} \emph{Résultat nommé}~: sur 200 ordres appariés, l'algorithme économise $-0.15$ point de base tous frais compris par rapport au TWAP, avec un intervalle à 95~\% de $[-0.44, 0.14]$~:
rien en moyenne~; l'économie vaut $1.25-0.70\log_2(Q/4000)-0.69\,h$ ($h=1$ en forte \omterm{def:mx:the-almgren-chriss-framework:urgency}{urgence}), de $+1.06$ sur des ordres de 4\,000 actions à faible \omterm{def:mx:the-almgren-chriss-framework:urgency}{urgence}
jusqu'à $-1.96$ sur 32\,000 en forte \omterm{def:mx:the-almgren-chriss-framework:urgency}{urgence}, le seuil de rentabilité se situant vers 13\,700 actions (environ 8~\% du volume).
\textbf{15.} Les exécutions passives économisent le spread et les frais proportionnellement aux actions~; l'impact croît avec la taille et l'\omterm{def:mx:the-almgren-chriss-framework:urgency}{urgence} et finit par les dépasser.
\textbf{16.} Environ 620 pour l'effet moyen, avec des nombres aléatoires communs~; les ordres réels ne partagent aucun contrefactuel, si bien que leur dispersion appariée est plus grande.
\textbf{17.} Un test pilote sur les petits ordres propres du desk, assorti d'un critère d'abandon, puis un test randomisé stratifié par taille, puis le réglage par défaut pour la strate gagnante.
\textbf{18.} Petits ordres à faible \omterm{def:mx:the-almgren-chriss-framework:urgency}{urgence}~; pour les gros ordres, une \omterm{def:mx:the-almgren-chriss-framework:urgency}{urgence} plus faible, une tolérance plus petite ou le TWAP.
\textbf{19.} En faire le réglage par défaut pour les ordres inférieurs à environ 8~\% du volume et garder le TWAP au-delà, en attendant un test en conditions réelles.
\textbf{20.} D'un point de base sur les petits ordres patients, et pas du tout sur les gros ordres urgents.
\end{solution}

\begin{solution}{iq:mx:build-an-execution-algorithm:1}
Il faut vérifier le motif (limite, bande, liquidité) et les instructions du client~; si la limite et l'\omterm{def:mx:the-almgren-chriss-framework:urgency}{urgence} le permettent, relever la participation dans la bande ou croiser
une partie du retard~; sinon, prévenir le client tôt que l'ordre risque de ne pas être exécuté en totalité, plutôt que de courir après le cours dans les dernières minutes.
\end{solution}

\begin{solution}{iq:mx:build-an-execution-algorithm:2}
Une fonction de décision pure de l'état (échéancier, exécutions, carnet, horloge), séparée de la messagerie et pilotée par une horloge simulée afin de pouvoir être
rejouée~; états et transitions explicites~; contrôles avant négociation regroupés en un seul endroit~; un journal d'audit de chaque transition, de chaque enfant et de chaque refus, avec son motif.
\end{solution}

\begin{solution}{iq:mx:build-an-execution-algorithm:3}
Randomiser les ordres entre les algorithmes (ou les apparier sur les mêmes ordres en simulation), mesurer le shortfall tous frais compris par rapport à l'arrivée, ajuster
selon la difficulté avec des erreurs groupées, et dimensionner le test à partir de la dispersion appariée~: $n=(z_{\alpha/2}+z_\beta)^2\sigma^2/\delta^2$.
\end{solution}

\begin{solution}{iq:mx:build-an-execution-algorithm:4}
Les contrôles propres de l'algorithme (taille des \omterm{def:mx:the-almgren-chriss-framework:parent}{ordres enfants}, collars, jamais au-delà de l'ordre, participation) attrapent ses erreurs~; la couche d'accès au marché du courtier
(seuils de crédit et de capital, contrôles des ordres erronés, coupe-circuit) doit se situer en dehors de lui, sous le contrôle direct de la firme, afin qu'aucun bogue de l'algorithme
puisse la contourner.
\end{solution}

\begin{solution}{iq:mx:build-an-execution-algorithm:5}
Que sa conception passive économise le spread et les frais sur les petits ordres, mais a davantage d'impact sur les gros~; il faut montrer l'attribution, la comparaison sur la
tranche de taille du client, et proposer une \omterm{def:mx:the-almgren-chriss-framework:urgency}{urgence} plus faible ou un autre algorithme pour ces ordres.
\end{solution}

\begin{solution}{iq:mx:build-an-execution-algorithm:6}
Les ordres n'ont pas été assignés au hasard~: les traders ont choisi les paramètres selon les conditions qu'ils observaient (l'\omterm{def:mx:the-almgren-chriss-framework:urgency}{urgence} les jours volatils), si bien que coûts et paramètres sont
confondus~; on y remédie par des variations aléatoires des paramètres sur une partie du flux, ou par un ajustement selon la difficulté et une simulation qui reproduit les choix.
\end{solution}
